\section{Analisi dei requisiti}

L'attività di analisi dei requisiti comprende la raccolta, la comprensione e la definizione delle esigenze del cliente e degli stakeholder. Il processo si svolge in tre fasi principali:
\begin{enumerate}
  \item \textbf{Raccolta delle informazioni}: interviste, documenti esistenti e casi d'uso per identificare gli obiettivi del sistema;
  \item \textbf{Organizzazione e validazione}: strutturare i requisiti in categorie funzionali e non funzionali, controllare la completezza e verificare con il cliente;
  \item \textbf{Formalizzazione}: redigere la documentazione ufficiale dei requisiti, includendo descrizioni chiare, criteri di accettazione e vincoli di progetto.
\end{enumerate}

Questa sezione stabilisce come si svolge il lavoro di analisi dei requisiti all'interno del processo di sviluppo, assicurando che ogni esigenza sia tracciabile e allineata agli obiettivi del progetto.
